\newif\ifPDF
\ifx\pdfoutput\undefined\PDFfalse
\else\ifnum\pdfoutput > 0\PDFtrue
	\else\PDFfalse
	\fi
\fi

\ifPDF
	\documentclass[pdftex,10pt]{article}
	\RequirePackage[hyperindex,colorlinks,plainpages=false]{hyperref}
	\hypersetup{pdfauthor={Heng Li},linkcolor=blue,citecolor=blue,urlcolor=blue}
	\usepackage{graphicx}
	\DeclareGraphicsRule{*}{mps}{*}{}
\else
	\documentclass[10pt]{article}
	\usepackage{graphicx}
\fi

\usepackage{amsmath}
\usepackage{amssymb}

\usepackage{natbib}
\bibliographystyle{apalike}

\usepackage{caption}
\renewcommand{\captionfont}{\fontsize{9pt}{11pt}\selectfont}
\newcommand{\lhspace}{\hspace{20pt}}

\addtolength{\textwidth}{3cm}
\addtolength{\hoffset}{-1.5cm}
\addtolength{\textheight}{4cm}
\addtolength{\voffset}{-2cm}

\title{TreeBeST}
\author{Heng Li and Richard Durbin}

\makeindex

\begin{document}

\section{Methods}

\subsection{Building gene tree with species-aware maximum-likelihood method}

\subsubsection{Construction of spec-SCFG}

Let $S$ be a binary rooted species tree, $G$ a binary rooted gene tree
with $n$ leaves and $X$ is the multiple alignment underlying
$G$. Assuming $X$ and $S$ independently influence a gene tree, we have
$P(G|S,X)=P(G|S)P(G|X)$. While probability $P(G|X)$ can be calculated
with Felsenstein's pruning algorithm, we will show in this section how
to approximate $P(G|S)$ with the help of the SCFG theory.

An SCFG generates a string instead of a tree. To calculate $P(G|S)$ we
need to write $G$ as a string. Given a gene tree $G$, let $x_G$ is
the New Hampshire string with a leaf of species $s$ written as
$g_s$. For example, for a gene tree $G$={\tt
  (gene1\_mouse,(gene2\_mouse,gene2\_rat)))}, a possible $x_G$ is
$(g_{mouse},(g_{mouse},g_{rat}))$. As swapping the branch order of an
internal node does not change the topology but may lead to a different
New Hampshire string, multiple $x_G$ may correspond to the same
topology. At the same time, not every swap of internal branch order will
change $x_G$. This is because the gene name has been omitted in
writting $x_G$, which brings additional symmetry. Let $m$ be the
number of swaps that do not change $x_G$. Then we have
$P(G|S)=2^{n-m-1}P(x_G|S)$. Probability $P(x_G|S)$ can be
calculated in the SCFG framework.

The SCFG can be constructed as follows. The non-terminal set is
$\{s|\mbox{$s$ is a node of $S$}\}$ The terminal set is
$\Sigma=\{g_s|\mbox{$s$ is a terminal node of $S$}\} \cup
\{\mbox{`(',`,',`)'}\}$ where $g_s$ represents a present gene of species
$s$. The rules are given in Table~\ref{tab:scfgrule}.

\begin{table}
\begin{center}
\begin{tabular}{l|lll}
\hline
Type & Rule & Probability & Meaning \\
\hline
$s$ is an  & $s\to(s_1,s_2)$ & $(1-p)(1-q)/2$ & internal speciation \\
internal   & $s\to(s_2,s_1)$ & $(1-p)(1-q)/2$ & internal speciation \\
node       & $s\to(s,s)$ & $p$ & duplication \\
           & $s\to s_1$ & $q(1-p)/2$ & loss \\
           & $s\to s_2$ & $q(1-p)/2$ & loss \\
\hline
$s$ is an  & $s\to g_s$ & $1-p$ & external speciation \\
external node & $s\to(s,s)$ & $p$ & duplication \\
\hline
\end{tabular}
\caption{Rules of species-SCFG. State $s$ represents the speciation
  state of a node of species tree $S$, $s_1$ and $s_2$ are the children
  of $s$ and $g_s$ is a gene of species $s$.}\label{tab:scfgrule}
\end{center}
\end{table}

\subsubsection{Probability of a gene tree}

Suppose we have a derivation consisting of $n_d$ duplication rules and
$n_l$ speciation rules. We can calculate this probability:
\begin{equation*}
  P(x_G,n,n_d,n_l|S)=\frac{p^{n_d}q^{n_l}}{2^{n-n_d+n_l-1}}(1-p)^{2n-n_d+n_l-1}(1-q)^{n-n_d-1}
\end{equation*}
and therefore:
\begin{equation}\label{equ:p-pi1}
  P(G,n,n_d,n_l|S)=2^{n_d-n_l-m}p^{n_d}q^{n_l}(1-p)^{2n-n_d+n_l-1}(1-q)^{n-n_d-1}
\end{equation}

Probability $P(x_G|S)$ can be calculated with the inside
algorithm. Given a node $g$ of $G$, let $x_g$ be the substring for the
subtree rooted at $g$, and $\alpha(x_g,s)$ be the probability that $x_g$
is generated from state $s$. If $g$ is a leaf node:
\begin{equation*}
  \alpha(x_g,s)=\left\{\begin{array}{ll} 1-p & \mbox{if $s$ is external and $g$ represents a gene of $s$} \\
    0 & \mbox{otherwise}\end{array}
  \right.
\end{equation*}
If $g$ is internal and its children are $g_1$ and $g_2$, $x_g$ can be
written as $x_g=(x_{g_1},x_{g_2})$. If $s$ is external:
\begin{equation*}
  \alpha(x_g,s)=\alpha((x_{g_1},x_{g_2}),s)=p\cdot\alpha(x_{g_1},s)\alpha(x_{g_2},s)
\end{equation*}
If both $s$ and $g$ are internal:
\begin{eqnarray*}
  \alpha(x_g,s)&=&p\cdot\alpha(x_{g_1},s)\alpha(x_{g_2},s)
  +\frac{q(1-p)}{2}\big[\alpha(x_g,s_1)+\alpha(x_g,s_2)\big]\\
  &&+\frac{(1-p)(1-q)}{2}\big[\alpha(x_{g_1},s_1)\alpha(x_{g_2},s_2)
  +\alpha(x_{g_1},s_2)\alpha(x_{g_2},s_1)\big]
\end{eqnarray*}
where $s_1$ and $s_2$ are the children of $s$. With the above three
equations we can calculate all $\alpha(x_g,s)$ and the probbility
$P(G|S)$ is:
\begin{equation}\label{equ:p-gs}
  P(G|S)=2^{n-m-1}P(x_G|S)=2^{n-m-1}\alpha(x_G,{\rm root}(S))
\end{equation}

\subsubsection{Practical implementation}

Although it is possible to calculate $P(G|S)$ as is showed in
Equation~\ref{equ:p-gs}, we did not implement this algorithm. First, in
real cases we frequently have multifurcated species tree due to lack of
evidence or to the short speciation time, but the SCFG described above
can only be appied to a binary species tree. Extention to multifurcated
species trees is hampered by the computational complexity when we intend
to traverse each possible binary topology at a multifurcated
node. Second, the inside algorithm is an $O(Mn)$ algorithm, where $M$ is
the size of the species tree and $n$ the size of the gene tree. When $M$
and $n$ are large, it is slower than Felsenstein's linear-time algorithm
for calculating $P(G|X)$. To avoid the complication of multifurcated
species tree and to achieve fast tree building, we decide to use an
approximate Equation~\ref{equ:p-pi1} to reflect the contribution of
species evolution.

In implementation, we use the following formula:
\begin{eqnarray*}
  P(G|S)&\approx&P(G,n,n^*_d,n^*_{l_s},n^*_{l_d}|S)\\
  &=&2^{n-m-1}P(G,n^*,n^*_d,n^*_{l_s},n^*_{l_d}|S)\\
  &=&2^{n^*_d-n^*_{l_s}-n^*_{l_d}-m}p^{n^*_d}q_s^{n^*_{l_s}}q_d^{n^*_{l_d}}
  (1-p)^{2n-n^*_d+n^*_{l_s}+n^*_{l_d}-1}(1-q_s)^{n-3n^*_d+n^*_{l_d}-1}(1-q_d)^{2n^*_d-n^*_{l_d}}\\
  &\approx&2^{n^*_d-n^*_{l_s}-n^*_{l_d}}p^{n^*_d}q_s^{n^*_{l_s}}q_d^{n^*_{l_d}}
  (1-p)^{2n-n^*_d+n^*_{l_s}+n^*_{l_d}-1}(1-q_s)^{n-3n^*_d+n^*_{l_d}-1}(1-q_d)^{2n^*_d-n^*_{l_d}}
\end{eqnarray*}
where $q_s$ is the probability of gene loss following a speciation and
$q_d$ the probability following a duplication. Number $n^*_d$ is the
minimum number of duplications, and $n^*_{l_s}$ and $n^*_{l_d}$ are the
minimum number of gene losses following a speciation or a duplication,
respectively. These numbers can be inferred from parsimonious tree
reconciliation. Again, $m$ is the number of New Hampshire string $x_G$
that correponds to $G$, with $0\leq m\leq n_d^*$. Calculating $m$ can be
done in $O(n^3)$ in the worst cases. It is also ommitted to achieve
efficiency.

The tree reconciliation algorithm is similar to (Zmasek and Eddy, 2001)
but with the extention to infer from a multifurcated species tree. The
algorithm is further improved to do duplication/loss inference and gene
tree rooting in $O(Mn)$ time, faster than $O(M^2n)$ of the original
implementation. Although the tree reconciliation algorithm has the same
time complexity as the inside algorithm, the former is faster due to
smaller constant in practice.

\subsection{Tree merge algorithm}

\subsubsection{Tree merge problem}

Let $V$ be a set of leaves, $\mathcal{G}=\{G_1,G_2,\ldots,G_N\}$ be a
set of {\it binary rooted} gene trees whose leaf set is $V$. We consider
two internal nodes are the same if and only if the sets of leaves
descending from the two nodes are identical. Then we can label all the
nodes in these trees as $\mathcal{U}(\mathcal{G})=\{1,2,\ldots,n\}$ and
define $\omega_i\subset V$ as the set of leaves that descend from node
$i$. For convenience, the root node, which will appear in each tree of
$\mathcal{G}$, is always labeled as $n$. We use $U(G)$ to denote the
node set of tree $G$, and therefore
$\mathcal{U}(\mathcal{G})=\bigcup_{G\in\mathcal{G}} U(G)$. We further
define:
\begin{equation*}
  \mathcal{U}_i(\mathcal{G})=\{k\in\mathcal{U}(\mathcal{G}):\omega_k\subset\omega_i\}
\end{equation*}
as the node set that is completely contained in $i$ and define:
\begin{equation*}
  \Omega_i(\mathcal{G})=\{G:\mbox{$G$ is a binary rooted tree with leaf set $\omega_i$ and $U(G)\subset\mathcal{U}_i(\mathcal{G})$}\}
\end{equation*}
Set $\Omega_n(\mathcal{G})$ consists of all possible topologies that are
consistent with $\mathcal{U}(\mathcal{G})$. It can be considered as the
tree space spanned by $\{G_1,\ldots,G_N\}$.  Obviously, all $G_i$,
$i\in[1,n]$, belong to $\Omega_n(\mathcal{G})$.

To describe the topologies in $\mathcal{G}$, we introduce an indicator
function $\delta_i(j,k)$ with $1\leq i,j,k\leq n$. The value of this
function equals to one if and only if: i) node $j$ and $k$ are $i$'s
children in a certain tree $G\in\mathcal{G}$ and ii) $j<k$; otherwise
$\delta_i(j,k)=0$. The second condition $j<k$ creates an ordering of the
children.

If we define a weight function $F(G)$ given a tree $G$, the general tree
merge problem is to find a tree $G^*\in\Omega_n(\mathcal{G})$ that
minimizes $F(G^*)$. For a general weight function, we have to traverse
all possible $G$ in $\Omega_n(\mathcal{G})$. This exhaustive search may
be formidable if the tree space is huge. However, finding the optimal
tree $G^*$ can be achieved in $O(N\cdot|V|^2)$ time if $F(G)$ can be
expressed as an additive sum over internal nodes:
\begin{equation}\label{equ:F}
  F(G)=\sum_{\begin{subarray}{1}l,j,k\in U(G)\\\delta_l(j,k)=1\end{subarray}}f(j,k)
\end{equation}
As nodes that cover the same leaf set have the same label, function
$f(j,k)$ is actually defined on two sets of leaves $\omega_i$ and
$\omega_j$, and independent of the topologies in both sets. This
is a strong requirement and we will come back to this point later.

\subsubsection{Tree merge algorithm}

For all leaf nodes, define $h^*(i)=0$; for an internal node $i$:
\begin{equation*}
  h^*(i)=\min_{G\in\Omega_i(\mathcal{G})}F(G)
  =\min_{G\in\Omega_i(\mathcal{G})}\sum_{\begin{subarray}{1}l,j,k\in U(G)\\\delta_l(j,k)=1\end{subarray}}f(j,k)
\end{equation*}
which equals to the weight of the best subtree up to node $i$. Function
$h^*(i)$ has the following recursive form:
\begin{equation}
  h^*(i)=\min_{\delta_i(j,k)=1}\big\{h^*(j)+f(j,k)+h^*(k)\big\}
\end{equation}
The calculation of $h^*(n)$ fulfills the procedure to find
the optimal tree $G^*$ in $\Omega_n(\mathcal{G})$. A more detailed
algorithm is given in Table~\ref{tab:merge}.

\begin{table}[!hbp]
\begin{center}
\begin{tabular}{|l|}
\hline
{\bf Function:}\\
\lhspace{\sc TreeMerge}$(G_1,G_2,\ldots,G_N)$\\
{\bf Input:} \\
\lhspace$N$ gene trees $G_1,G_2,\ldots,G_N$ which have the same leaf set; \\
{\bf Output:} \\
\lhspace $U(G^*)$: the node sets of the optimal tree $G^*$\\
{\bf Procedures:} \\
\lhspace$selected\gets\mbox{empty array}$\\
\lhspace{\sc OptimizeF}$(i)$ \\
\lhspace\lhspace{\bf if} $h^*(i)$ has been calculated {\bf then} \\
\lhspace\lhspace\lhspace{\bf return} $h^*(i)$ \\
\lhspace\lhspace$min\leftarrow\infty$ \\
\lhspace\lhspace{\bf for} each $(j,k)$ satisfying $\delta_i(j,k)=1$ {\bf do} \\
\lhspace\lhspace\lhspace$score\gets\mbox{{\sc OptimizeF}}(j)+f(j,k)+\mbox{{\sc OptimizeF}}(k)$\\
\lhspace\lhspace\lhspace{\bf if} $min>score$ {\bf then}\\
\lhspace\lhspace\lhspace\lhspace$min\gets score$\\
\lhspace\lhspace\lhspace\lhspace$minpair\gets (j,k)$\\
\lhspace\lhspace$h^*(i)\gets min$ \\
\lhspace\lhspace$selected[i]\gets minpair$ \\
\lhspace\lhspace{\bf return} $h^*(i)$ \\
\\
\lhspace{\sc GetTreeNodes}$(i)$ \\
\lhspace\lhspace{\bf if} $i$ is a leaf {\bf then}\\
\lhspace\lhspace\lhspace{\bf return} $\{i\}$\\
\lhspace\lhspace$(j,k)\gets selected[i]$\\
\lhspace\lhspace{\bf return} $\mbox{\sc GetTreeNodes}(j)\cup\{i\}\cup\mbox{\sc GetTreeNodes}(k)$\\
\\
\lhspace{\sc TreeMerge}$(G_1,G_2,\ldots,G_N)$\\
\lhspace\lhspace Calculate $\mathcal{U}(\mathcal{G})$ and all $\delta_i(j,k)$\\
\lhspace\lhspace{\bf for} each $i$ is a leaf node {\bf do}\\
\lhspace\lhspace\lhspace$h^*(i)\gets 0$\\
\lhspace\lhspace{\sc OptimizeF}$(n)$\\
\lhspace\lhspace{\bf return} {\sc GetTreeNodes}$(n)$\\
\hline
\end{tabular}
\caption[Tree merge algorithm]{Tree merge algorithm. Function {\sc
    TreeMerge}() gives the node set of the optimal tree
  $G^*\in\Omega_n(\mathcal{G})$. The topology of $G^*$ could be
  recovered from this set.}~\label{tab:merge}
\end{center}
\end{table}

\subsubsection{Defining weight function}

As is described above, weight function $f(j,k)$ is defined on $\omega_j$
and $\omega_k$ only. The likelihood function used in Felsenstein's
algorithm and the parsimonious function do not satisfy this
requirement. For $\delta_i(j,k)=1$, a possible form of $f(j,k)$ is:
\begin{equation*}
  f(j,k)=c_df_d(j,k) + c_lf_l(j,k) - c_bf_b(j,k)
\end{equation*}
where $f_d(j,k)$ is an indicator function whether $i$ is a duplication;
$f_l(j,k)$ is the number of losses that happen between $(i,j)$ and
$(i,k)$; $f_b(j,k)$ is the bootstrap value on the branch between $i$ and
its parent; $c_d$, $c_l$ and $c_b$ are positive constants.

\end{document}